\subsection{Archive-consistent certification of saved ML forecasts}
\label{sec:historical-application}
The historical study uses 30 Portugal consumption series, one French
household series and three recorded cumulative-counter series. The four
forecasters are persistence, ridge regression, histogram gradient boosting
and a multilayer perceptron. Forecasts, chronological training periods,
eligibility filters and masks are reused unchanged. Appendix~\ref{app:historical-design}
specifies the protocol, training history and information boundary.

The archive-aware evaluator receives all saved test forecast rows within
the declared period. In particular, a saved persistence forecast
$p_{t+24}^{\rm pers}=y_t$ reveals its earlier target directly. These copies
and exact aggregate deductions are applied before all methods, using
source-row traces; hidden reference targets only check the recovered
values. No inverse reconstruction from other fitted predictor mappings is
assumed. Later records within the evaluation period are included. The
whole-period feasible set is the product of the remaining daily blocks
after this closure, and equal-length daily endpoints may therefore be
averaged. This is a retrospective benchmark archive rule, not a measured
production arrival deadline.

The complete Portugal archive-aware exact-total grid contains 9,720
repeated model-pair comparisons across gap lengths, mechanisms, seeds and
aggregate resolutions. Direct paired bounds certify 7,836 (80.62\%),
whereas independently bounded individual losses certify 2,513 (25.85\%).
The additional 5,323 certificates are a 54.76-percentage-point difference
under identical information and zero additional queries. In the six-hour-gap
illustrative cell the counts are 478 of 540 versus 64 of 540. The broader
grid, rather than that cell, supplies the abstract's empirical result.

Independent subtraction of sharp individual-loss bounds encloses the
paired interval but may combine extremizers that cannot be the same
record. For example, predictions $(1,1)$ and $(2,2)$ under
$y_1+y_2=2$, $y\geq0$, have individual MSE ranges $[0,1]$ and $[1,2]$;
separate subtraction gives $[-2,0]$, whereas the paired contrast is exactly
$-1$. For exact totals, each individual-loss minimum is a simplex
projection and each maximum is attained at a simplex vertex. For uncertain
totals, the remaining block instead satisfies
$x\geq0$ and $L\leq\sum_i x_i\leq U$, after subtracting observed contributions.
The minimum is the projection of its prediction vector onto this band.
If the nonnegative part of that vector has a sum in $[L,U]$, it is already
optimal and the minimizing total may be interior. Otherwise the projection
uses the lower or upper boundary simplex, respectively. The maximum of
the convex individual squared loss is attained at a vertex of the compact
band polytope.

The complete 34-series scenario grid contains 66,096 repeated comparisons:
10,404 are certified by both methods, 23,906 only by pairing and 31,786
by neither. None is certified only by separate-loss bounds. All 324
scenario cells, including unresolved cases, appear in the supplement;
Appendix~\ref{app:certificate-grid} displays their 18 cohort/access/tolerance
aggregates. The checker verifies the saved sign
rule, paired-interval nesting and reference enclosure. The redistributable
endpoint replay covers 60,264 Portugal/household comparisons; the 5,832
counter comparisons use local original endpoints and require the documented
raw-source regeneration for an external replay. These are different
reproduction levels, not interchangeable counts of independent evidence.

The audit adds all-completions support to a point evaluation. The historical
donor makes zero observed
whole-period errors at full coverage on the 9,720 Portugal, 324 household
and 972 counter exact-total comparisons. Those empirical successes do not
give it an all-completions guarantee. Our certificate can accompany a point
evaluation by identifying which rankings its declared information actually
forces; this study does not measure consequential downstream decisions
changed by that certificate.

\subsection{The retained adverse recovery result}
On the original Portugal six-hour-gap archive, saved persistence copies
reconstruct 95.33\% of masked target occurrences for free. Daily certification
already reaches 96.67\% before paid recovery and 98.71\% after at most two
readings per model pair; the exact-over-greedy addition is 0.042 percentage
points. A budget for one pair is not a shared budget for all six pairs.
No monetary retrieval cost or field latency was measured.

In the archived randomized comparison, all policies see the same closed
archive. On Portugal, normalized worst width decreases from 29.595\% to
26.673\% of initial width, but expected width at the actual record increases
from 20.154\% to 23.670\% (17.45\% relative). Certification decreases
from 98.795\% to 98.688\%. These are exact expectations over the computed
mixtures, not Monte Carlo estimates. Width fractions first sum within
each meter and then average meter ratios. The two-reading budget applies
separately to each day and pair, including already certified comparisons.

Across 1,074 positive-width public cases, expected realized width improves
in 14, worsens in 649 and ties in 411. Both target-independent refinements
in Section~\ref{sec:tie-selection} completed on all 1,074 cases without
removing the penalty beyond the declared tolerance. This does not imply
every optimal mixture has identical realized performance. The new theorem
establishes that adaptation cannot improve the robust mean-width value;
it does not turn these adverse observed outcomes into an operational benefit.
The compact solver is therefore an exact risk optimizer, not an empirically
dominant default acquisition rule.

